\documentclass[11pt]{article}
\usepackage{cs152}

%%%%%%%%%%%%%%%%%%%% name/id
\rfoot{\small Brian Park | 3033045069}


%%%%%%%%%%%%%%%%%%%% Course/HW info
\newcommand*{\instr}{John Wawrzynek}
\newcommand*{\term}{Spring 2022}
\newcommand*{\coursenum}{CS 152}
\newcommand*{\coursename}{Computer Architecture and Engineering}
\newcommand*{\hwnum}{3}

%%%%%%%%%%%%%%%%%%%%%%%%%%%%%% Document Start %%%%%%%%%%%%%%%%%
\begin{document}

%%%%%%%%%%%%%%%%%%%%%%%%%%%%%%%%%%%%%%%%%%%%%%%%%%%%%%%%%%%%%%%%%%%%%%%%%%%%%%%%%%%%%%%%
% Question 3
%%%%%%%%%%%%%%%%%%%%%%%%%%%%%%%%%%%%%%%%%%%%%%%%%%%%%%%%%%%%%%%%%%%%%%%%%%%%%%%%%%%%%%%%
\section*{(Directed) Problem 3.3: Performance Bottlenecks}



\begin{table}[ht]
\caption{CPI for Different Optimizations} % title of Table
\centering % used for centering table
\begin{tabular}{c c c c c c c c} % centered columns (4 columns)
\hline\hline %inserts double horizontal lines
& dhrystone & median & multiply & qsort & spmv & towers & vvadd \\ [0.5ex] % inserts table
%heading
\hline % inserts single horizontal line
Rocket & 1.200 & 1.379 & 1.136 & 1.427 & 1.690 & 1.029 & 1.024 \\
BOOM (worst-case) & 6.679 & 5.538 & 7.275 & 5.698 & 6.001 & 6.144 & 6.884 \\
BOOM (64 PRF) & 2.891 & 3.674 & 3.743 & 3.126 & 2.627 & 2.197 & 2.779 \\
BOOM (16 ROB) & 6.671 & 5.446 & 7.275 & 5.698 & 5.906 & 6.103 & 6.884 \\
BOOM (br. pred.) & 5.886 & 4.8561 & 6.568 & 5.510 & 5.905 & 5.876 & 6.503 \\
[1ex] % [1ex] adds vertical space
\hline %inserts single line
\end{tabular}
\label{table:nonlin} % is used to refer this table in the text
\end{table}

\begin{table}[ht]
\caption{Branch Misprediction Accuracies} % title of Table
\centering % used for centering table
\begin{tabular}{c c c c c c c c} % centered columns (4 columns)
\hline\hline %inserts double horizontal lines
& dhrystone & median & multiply & qsort & spmv & towers & vvadd \\ [0.5ex] % inserts table
%heading
\hline % inserts single horizontal line
%Rocket & 1.200 & 1.379 & 1.136 & 1.427 & 1.690 & 1.029 & 1.024 \\
BOOM (worst-case) & 0.589 & 0.570 & 0.904 & 0.611 & 0.476 & 0.827 & 0.766 \\
BOOM (64 PRF) & 0.589 & 0.570 & 0.904 & 0.611 & 0.476 & 0.827 & 0.766 \\
BOOM (16 ROB) & 0.589 & 0.570 & 0.904 & 0.611 & 0.476 & 0.827 & 0.766 \\
BOOM (br. pred.) & 0.002 & 0.173 & 0.087 & 0.177 & 0.192 & 0.070 & 0.032 \\
[1ex] % [1ex] adds vertical space
\hline %inserts single line
\end{tabular}
\label{table:nonlin} % is used to refer this table in the text
\end{table}



\begin{Parts}
	\Part Compare the results for 1-wide BOOM with the in-order, 5-stage Rocket core. Was out-of-order issue an improvement on the CPI for these benchmarks? Explain the impact of each microarchitectural change. \\
	Yes, in general, the out of order issue in BOOM was an improvement for each microarchitectural change. There is always some extent of improvement of CPI in all benchmarks and all microarchitectural changes. For just BOOM in the worst case, it exemplifies the best improvement compared to the other microarchitectural changes. The OoO execution definitely helps in exploiting ILP over the Rocket processors, but as more hardware complexity is added, the CPI performance dips down again. We see it dip down the most for increasing the physical register file to 64. This may be due to the increased latency to index the registers. The most minimal performance dip is changing the reorder buffer from 4 to 16. This minimal dip might be because of the increased overhead of holding more entries in the reorder buffer. And for enabling branch prediction, we see another dip in CPI. This dip could be due to more branch mispredicts and the cost it incurs to take the correct branch again. 
	\Part Is branch prediction with a BTB, BHT, and RAS always a benefit? \\
	In this case, yes. Branch prediction across all benchmarks provide better prediction accuracies. For even smaller programs though, the costs may outweigh the benefits, where the program will not have enough time to be trained on the program's branch pattern. Also, since hardware complexity is added, it may increase the space used on the chip dye, affecting engineering costs.
\end{Parts}
\newpage

%%%%%%%%%%%%%%%%%%%%%%%%%%%%%%%%%%%%%%%%%%%%%%%%%%%%%%%%%%%%%%%%%%%%%%%%%%%%%%%%%%%%%%%%
% Question 4
%%%%%%%%%%%%%%%%%%%%%%%%%%%%%%%%%%%%%%%%%%%%%%%%%%%%%%%%%%%%%%%%%%%%%%%%%%%%%%%%%%%%%%%%
\section*{(Open-ended) Problem 4.1: Designing Your Own Branch Predictor}
\section*{(Open-ended) Problem 4.2: Recreating Spectre Attacks}
\newpage

%%%%%%%%%%%%%%%%%%%%%%%%%%%%%%%%%%%%%%%%%%%%%%%%%%%%%%%%%%%%%%%%%%%%%%%%%%%%%%%%%%%%%%%%
% Question 5
%%%%%%%%%%%%%%%%%%%%%%%%%%%%%%%%%%%%%%%%%%%%%%%%%%%%%%%%%%%%%%%%%%%%%%%%%%%%%%%%%%%%%%%%
\section*{(Directed) Problem 5: Feedback}
The directed portion took around 10 hours. Mostly, it was getting setup and running scripts, which was done asynchronously for me while doing other stuff. This lab took really long with compilation/simulation. For the future, I wish there are options to run locally on a Docker or VM. I would like to take advantage of my own hardware if that speeds up anything.

We spent nearly 10 hours on open-ended portion as well. This was really fun reverse engineering. Usually, I use Linux tools like \verb|htop| or \verb|lscpu| to know my hardware. This time it was fun to blindly figure out by reverse engineering. Can't wait to see if our results were at all close to the mystery configuration!

I learned how to empirically analyze and engineer assembly/C code. I think having this intuition and knowledge behind performance is definitely beneficial in my future career and current works in research!

There is nothing that I would like to change.

Contribution was done individually for each directed part. Then we shared work evenly for open ended, each trying different approaches. I went to OH and found some of the input from other classmates and Josh helpful. 
My partner Kevin Chang contributed mostly to finding more about the data cache, while I contributed by reverse engineering specifications about the instruction cache.

\end{document}
